\section*{الفصل \ref{ch:b2:oceanic-islands} --- الجزر المحيطية}

\begin{solution}{exo:b2:oceanic-islands:1}
(أ) مرحلة ما قبل الدرع، مثل لوئيهي اليوم. (ب) \omterm{def:b2:oceanic-islands:stages}{مرحلة الدرع}. (ج) مرحلة ما بعد الدرع. (د) مرحلة
التجدّد. (ه) لم تعد أي مرحلة ثورانية: فقد غاص البركان تحت شعابه، التي صارت الآن حلقة مرجانية.
\end{solution}

\begin{solution}{exo:b2:oceanic-islands:2}
(أ) $2500 + 350\sqrt9 = \qty{3550}{m}$. (ب) $350\times(\sqrt{25} - \sqrt9) = 350\times2 =
\qty{700}{m}$. (ج) $350/(2\times5) = \qty{35}{m}$ لكل مليون سنة.
\end{solution}

\begin{solution}{exo:b2:oceanic-islands:3}
(أ) $5000/17 = \qty{294}{kyr}$ حول الجزر الرئيسية؛ و$30\,000/70 = \qty{429}{kyr}$ على امتداد
الظهرة. (ب) تفوت بعض الانزلاقات: فيمكن لبركان أحدث أن يدفن ندبة بركان أقدم وحطامه،
والرواسب القديمة تغطّيها الرسوبيات. فالعدد الحقيقي أكبر والفترة الحقيقية
أقصر.
\end{solution}

\begin{solution}{exo:b2:oceanic-islands:4}
(أ) $4158\times0.3048 = \qty{1267}{m}$؛ و$4610\times0.3048 = \qty{1405}{m}$؛ و$150\times0.3048 =
\qty{46}{m}$. (ب) نمت أولى الشعاب على عمق لا يزيد على \qty{46}{m} وقاعدتها اليوم عند
\qty{1267}{m}: فقد غاص البركان بما لا يقلّ عن $1267 - 46 = \qty{1221}{m}$ (وأكثر إن كان مستوى البحر
أعلى حينئذ). (ج) الغوص نحو \qty{1267}{m} على الأكثر، خلال \qty{33.9}{Myr} على الأقلّ:
أي $1267/33.9 = \qty{37}{m}$ لكل مليون سنة على الأكثر، أو \qty{0.037}{mm} في السنة.
\end{solution}

\begin{solution}{exo:b2:oceanic-islands:5}
(أ) $\dot S = 175/\sqrt t$: \qty{87.5}{m.Myr^{-1}} (\qty{0.0875}{mm.yr^{-1}}) عند \qty{4}{Myr}،
و\qty{35}{m.Myr^{-1}} (\qty{0.035}{mm.yr^{-1}}) عند \qty{25}{Myr}، و\qty{25}{m.Myr^{-1}}
(\qty{0.025}{mm.yr^{-1}}) عند \qty{49}{Myr}. (ب) $4.2/0.0875 = 48$ مرة. (ج) $2.6/0.0875 = 30$
مرة: فمرحلة التحميل تُغوِّص الجزيرة أسرع بكثير من تبرّد صفيحة فتيّة.
\end{solution}

\begin{solution}{exo:b2:oceanic-islands:6}
(أ) $\dot S = (3200/62.8)\,e^{-100/62.8} = 50.96\times0.2035 = \qty{10.37}{m}$ لكل مليون سنة،
أي \qty{0.0104}{m} لكل ألف سنة: وهو المعدّل المقيس \qty{0.01}{m} لكل ألف سنة. (ب) $e^{-t/62.8}
= 10/50.96$، و$t = 62.8\ln5.096 = \qty{102.3}{Myr}$. (ج) قانون الجذر: $2500 + 350\times5 =
\qty{4250}{m}$؛ والقانون الأسّي: $6400 - 3200\times0.6716 = \qty{4251}{m}$؛ والفرق بينهما
\qty{1}{m}.
\end{solution}

\begin{solution}{exo:b2:oceanic-islands:7}
(أ) $\qty{8.6}{cm.yr^{-1}} = \qty{86}{km.Myr^{-1}}$؛ و$86\times33 = \qty{2838}{km}$. (ب) نعم:
\qty{27.7}{Myr} أقلّ من \qty{33}{Myr} التي تغرق الشعاب بعدها، فينبغي أن تكون شعابها
لا تزال حيّة. (ج) $55\times1.609 = \qty{88.5}{km}$، أي $88.5/86 = \qty{1.0}{Myr}$ إضافية: نحو
\qty{28.7}{Myr}. وهذا أقلّ إلى حدّ ما من \qty{33}{Myr}، لكن قيمة \qty{33}{Myr} نفسها رقم
تقريبي للسلسلة كلّها، والنموّ يتوقّف أيضًا على حرارة الماء عند كل
حلقة مرجانية: فكوري، آخر حلقة مرجانية حيّة، قريبة فعلًا من \omterm{def:b2:oceanic-islands:darwin-point}{نقطة داروين}.
\end{solution}

\begin{solution}{exo:b2:oceanic-islands:8}
(أ) الصندوق: $30\times17\times2 = \qty{1020}{km^3}$؛ والهرم: الثلث، \qty{340}{km^3}. (ب) كتلة واحدة
وحدها عُشر (الصندوق) أو جزء من ثلاثين (الهرم) من بركان كيلاويا كلّه. (ج) $23\,000\times
0.1 = \qty{2300}{km^3}$، أي ربع كيلاويا؛ وهو تقدير يصحّ أو يسقط مع السمك
المفترض.
\end{solution}

\begin{solution}{exo:b2:oceanic-islands:9}
(أ) $1000/2.6 = \qty{385}{kyr}$. (ب) نعم: \qty{2.6}{m} لكل ألف سنة تساوي
\qty{2.6}{mm.yr^{-1}}، وهي أقلّ من \qty{4.2}{mm.yr^{-1}}، بهامش \qty{1.6}{mm.yr^{-1}}.
(ج) الارتفاع النسبي $12 + 2.6 = \qty{14.6}{mm.yr^{-1}}$؛ فتتخلّف الشعاب بمقدار $14.6 - 4.2 =
\qty{10.4}{mm.yr^{-1}}$ خلال $16.5 - 8.2 = \qty{8.3}{kyr}$: فتنتهي قمّتها عند $10.4\times8.3 =
\qty{86}{m}$. (د) $(12 - 4.2)\times8.3 = \qty{65}{m}$. وفي الحالتين تبقى فوق أعمق
حدّ للنموّ لكن في ماء أعتم بكثير: أمّا هل تلحق بالسطح فيتوقّف على كيف يتباطأ نموّها مع
العمق.
\end{solution}

\begin{solution}{exo:b2:oceanic-islands:10}
(أ) للدالة $S(t_1) = C(\sqrt{t_1 + \Delta} - \sqrt{t_1})$ المشتقّ $\frac C2(1/\sqrt{t_1 +
\Delta} - 1/\sqrt{t_1}) < 0$: فهي تتناقص مع $t_1$. (ب) $350\sqrt{20} = \qty{1565}{m}$؛ و$350\times
(\sqrt{70} - \sqrt{50}) = 350\times1.2955 = \qty{453}{m}$. (ج) لا: $\sqrt t$ منتهٍ عند $0$، لذا
فالغوص على أي فترة منتهٍ، \qty{1565}{m} في (ب). ولا يكون لانهائيًا إلا المعدّل في اللحظة $t = 0$،
حيث لم يعد القانون معدًّا لوصف القاع بالتفصيل.
\end{solution}

\begin{solution}{exo:b2:oceanic-islands:11}
(أ) $30\times0.1 = \qty{3}{m}$. (ب) نحو \qty{120}{m} إن كانت قمّتها عند مستوى البحر في
الارتفاع الأعلى: فتصير \omterm{def:b2:oceanic-islands:reefs}{الحلقة المرجانية} جزيرة من الحجر الجيري، ويذيب المطر أراغونيت المرجانيات
ويُعاد تبلور الصخر إلى كالسيت، مع كهوف. (ج) $120/3 = 40$ دورة، أي \qty{4}{Myr}. (د) النطاقات
المغسولة هي أزمنة الانكشاف هذه، مختومة تحت الشعاب التي نمت من جديد حين عاد البحر.
وعلى مدى دورات كثيرة يعود البحر إلى ارتفاعات عليا متشابهة، فيكون $\Delta\eta$ صغيرًا، وتعود قمّة الشعاب
فتنمو حتى مستوى البحر عند كل ارتفاع أعلى، ويحدّد الغوصُ سمك الغطاء، كما تقول
القضيّة.
\end{solution}

\begin{solution}{exo:b2:oceanic-islands:12}
(أ) $3200\times(e^{-50/62.8} - e^{-120/62.8}) = 3200\times(0.4510 - 0.1480) = \qty{970}{m}$.
(ب) $150 + 970 = \qty{1120}{m}$. (ج) $700\times1.8288 = \qty{1280}{m}$، أي أعمق من
التقدير بمقدار \qty{160}{m}. (د) ربما غرق \omterm{def:b2:oceanic-islands:seamount}{الغيوت} أبكر، على صفيحة أفتى، فغاص أكثر؛ وربما حمّله بركان
مجاور؛ وربما قطعت الأمواج قمّته أخفض، أو تغطّيها رسوبيات؛ وقد تكون
الصفيحة تحته من عمر مختلف؛ ولم يكن مستوى البحر في زمن الغرق هو مستوى
اليوم.
\end{solution}

\begin{solution}{pb:b2:oceanic-islands:1}
\textbf{1.} $2500 + 350\sqrt{10} = 2500 + 350\times3.1623 = \qty{3607}{m}$.
\textbf{2.} $3607 + 1000 = \qty{4607}{m}$.
\textbf{3.} $V = \pi r^2h/3 = \pi\times50^2\times4.607/3 = \qty{12061}{km^3}$.
\textbf{4.} $0.85\times12\,061 = \qty{10252}{km^3}$ في \qty{1}{Myr}: أي \qty{0.0103}{km^3} في السنة،
أقلّ بعشر مرات من كيلاويا اليوم.
\textbf{5.} $95\,600/12\,061 = 7.9$: نحو ثمانية.
\textbf{6.} الحمم الثولييتية، من \omterm{def:b2:oceanic-islands:stages}{مرحلة الدرع}؛ والحمم القلوية لمرحلة ما بعد الدرع
تغطّيه.
\textbf{7.} $\qty{1}{km}/\qty{1}{Myr} = \qty{1}{m}$ لكل ألف سنة، أبطأ بمقدار 2.6 مرة مما قبالة
جزيرة هاواي.
\textbf{8.} نعم: $4.2 - 1.0 = \qty{3.2}{mm.yr^{-1}}$ فائضة.
\textbf{9.} $12 + 1 - 4.2 = \qty{8.8}{mm}$ في السنة.
\textbf{10.} $150/8.8 = \qty{17}{kyr}$. مصطبة شعاب غارقة على الجناح، درجة من النوع
الموجود قبالة جزيرة هاواي.
\textbf{11.} $350\times(\sqrt{22} - \sqrt{10}) = 350\times(4.6904 - 3.1623) = \qty{535}{m}$.
\textbf{12.} $1000 - 535 = \qty{465}{m}$، بمعدّل $465/12 = \qty{38.8}{m}$ لكل مليون سنة.
\textbf{13.} $350/(2\sqrt{10}) = \qty{55.3}{m.Myr^{-1}}$ في البداية؛ و$350/(2\sqrt{22}) =
\qty{37.3}{m.Myr^{-1}}$ في النهاية.
\textbf{14.} $2500 + 350\times4.6904 = \qty{4142}{m}$.
\textbf{15.} تعلو جزر هاواي الكبيرة حتى نحو أربعة أضعاف ذلك، فيكون فوق مستوى البحر صخر أكثر بكثير
يجب إزالته؛ كما دامت \omterm{def:b2:oceanic-islands:stages}{مرحلة الدرع} فيها أطول.
\textbf{16.} الأمواج، التي تُرجِع الساحل وتسوّي مصطبة نحت الأمواج؛ والانزلاقات وانهيارات
الأجنحة.
\textbf{17.} شعاب حاجزية، بينها وبين آخر ما بقي من الشاطئ البركاني بحيرة.
\textbf{18.} $350\times(\sqrt{70} - \sqrt{22}) = 350\times(8.3666 - 4.6904) = \qty{1287}{m}$.
\textbf{19.} $1287/48 = \qty{26.8}{m}$ لكل مليون سنة، أي \qty{0.027}{mm} في السنة: $4.2/0.0268 =
157$ مرة أقلّ مما كان يمكن للشعاب أن تنمو.
\textbf{20.} أسمك بمقدار \qty{20}{m} من الحجر الجيري فوق البازلت تحت جزيرة باري
(\qty{1267}{m})، وأرقّ بمقدار \qty{118}{m} مما تحت إلوغيلاب (\qty{1405}{m}): فهما من الرتبة نفسها.
\textbf{21.} $2500 + 350\times8.3666 = \qty{5428}{m}$.
\textbf{22.} $3200\times(e^{-70/62.8} - e^{-120/62.8}) = 3200\times(0.3280 - 0.1480) =
\qty{576}{m}$.
\textbf{23.} $150 + 576 = \qty{726}{m}$.
\textbf{24.} يبلغ سمك غطاء الشعاب لحلقة بترل المرجانية \textbf{\qty{1287}{m}}، أي نحو \qty{1.3}{km}
من الحجر الجيري الشعابي، كما تحت إنيويتاك.
\end{solution}
